% ============================================================
\chapter{Lagrangian Duality}
\label{ch:lagrangian}
% ============================================================

Lagrangian duality is one of the most unifying ideas in optimisation. Its principle: relax the constraints by penalising them in the objective via \emph{Lagrange multipliers}, then maximise over these multipliers. The resulting dual problem is always convex (even if the primal is not), provides a guaranteed lower bound, and coincides with the primal under \emph{strong duality} conditions---which are automatically satisfied in convex optimisation as soon as a Slater constraint qualification holds. This theory, heir to the work of Lagrange (1797), Kuhn-Tucker (1951), and Rockafellar (1970), is the bridge between the primal formulation and the KKT conditions.

\section{Introduction}

Lagrangian duality is a fundamental principle that associates to every optimization
problem (primal) a dual problem. The dual provides lower bounds on the primal optimal
value and, under convexity and qualification conditions, strong duality allows solving
the primal via the dual.

\section{The Lagrangian}

\begin{definition}[Lagrangian]
The \emph{Lagrangian} of problem (P) is the function
$L : \R^n \times \R^m_+ \times \R^p \to \R$ defined by:
\[
    L(x, \lambda, \nu) = f(x) + \sum_{i=1}^m \lambda_i g_i(x) + \sum_{j=1}^p \nu_j h_j(x).
\]
The variables $\lambda \in \R^m_+$ and $\nu \in \R^p$ are the \emph{dual variables}
or \emph{Lagrange multipliers}.
\end{definition}

\begin{proposition}[The Lagrangian as Relaxation]
For all $\lambda \geq 0$ and $\nu$, and all feasible $x$:
\[
    L(x, \lambda, \nu) \leq f(x).
\]
Indeed, $\lambda_i g_i(x) \leq 0$ and $\nu_j h_j(x) = 0$.
\end{proposition}

\section{Lagrange Dual Function}

\begin{definition}[Dual Function]
The \emph{Lagrange dual function} is:
\[
    q(\lambda, \nu) = \inf_{x \in \R^n} L(x, \lambda, \nu)
    = \inf_x \left\{f(x) + \sum_i \lambda_i g_i(x) + \sum_j \nu_j h_j(x)\right\}.
\]
\end{definition}

\begin{theorem}[Properties of the Dual Function]
\label{thm:dual}
\begin{enumerate}
    \item $q$ is concave (as the infimum of affine functions in $(\lambda, \nu)$).
    \item \textbf{Weak duality}: $q(\lambda, \nu) \leq p^*$ for all $\lambda \geq 0$
    and $\nu$, where $p^*$ is the primal optimal value.
    \item The domain $\mathrm{dom}(q) = \{(\lambda,\nu) \mid q(\lambda,\nu) > -\infty\}$
    is convex.
\end{enumerate}
\end{theorem}

\begin{proof}[Proof of weak duality]
For any feasible $x$ and any $(\lambda, \nu)$ with $\lambda \geq 0$:
\[
    q(\lambda, \nu) = \inf_{z} L(z, \lambda, \nu) \leq L(x, \lambda, \nu)
    = f(x) + \sum_i \underbrace{\lambda_i g_i(x)}_{\leq 0}
    + \sum_j \underbrace{\nu_j h_j(x)}_{= 0} \leq f(x).
\]
Taking the infimum over feasible $x$: $q(\lambda, \nu) \leq p^*$.
\end{proof}

\section{Dual Problem}

\begin{definition}[Lagrange Dual Problem]
The \emph{dual problem} is:
\[
    (D) \quad d^* = \sup_{\lambda \geq 0, \, \nu} q(\lambda, \nu)
    = \sup_{\lambda \geq 0, \, \nu} \inf_x L(x, \lambda, \nu).
\]
\end{definition}

\begin{definition}[Duality Gap]
The \emph{duality gap} is $p^* - d^* \geq 0$. If $p^* = d^*$, we say that
\emph{strong duality} holds.
\end{definition}

\begin{keyformulas}[title=\textbf{Primal--Dual Relationship}]
\[
    d^* = \sup_{\lambda \geq 0, \nu} \inf_x L(x,\lambda,\nu)
    \leq \inf_x \sup_{\lambda \geq 0, \nu} L(x,\lambda,\nu) = p^*.
\]
The inequality is the max-min inequality.
\end{keyformulas}

\section{Strong Duality}

\begin{theorem}[Strong Duality --- Slater]
\label{thm:slater}
If the primal problem is convex ($f, g_i$ convex, $h_j$ affine), $p^*$ is finite,
and Slater's condition is satisfied, then:
\begin{enumerate}
    \item Strong duality holds: $p^* = d^*$.
    \item The dual supremum is attained: there exist optimal $(\lambda^*, \nu^*)$.
\end{enumerate}
\end{theorem}

\begin{proof}[Proof sketch]
Consider the set:
\[
    \mathcal{G} = \{(u, v, t) \in \R^m \times \R^p \times \R \mid
    \exists x : g_i(x) \leq u_i, h_j(x) = v_j, f(x) \leq t\}.
\]
$\mathcal{G}$ is convex (by convexity of $f, g_i$ and affinity of $h_j$).
The point $(0, 0, p^*)$ lies on the boundary of $\mathcal{G}$. By the supporting
hyperplane theorem, there exists a separating hyperplane, which yields the optimal
Lagrange multipliers.
\end{proof}

\section{Examples of Dual Computation}

\begin{example}[Linear Programming]
Primal: $\min c^Tx$ subject to $Ax \leq b$, $x \geq 0$.

$L(x, \lambda, \mu) = c^Tx + \lambda^T(Ax - b) - \mu^Tx = (c + A^T\lambda - \mu)^Tx - \lambda^Tb$.

$q(\lambda, \mu) = \inf_x L = \begin{cases} -\lambda^Tb & \text{if } c + A^T\lambda - \mu = 0, \\ -\infty & \text{otherwise.}\end{cases}$

Dual: $\max -b^T\lambda$ subject to $A^T\lambda \leq c$, $\lambda \geq 0$.

After standard reformulation: the dual of an LP is an LP.
\end{example}

\begin{example}[Minimum Norm]
Primal: $\min \norm{x}^2$ subject to $Ax = b$.

$L(x, \nu) = \norm{x}^2 + \nu^T(Ax - b)$.

Optimality condition $2x + A^T\nu = 0$ gives $x = -\frac{1}{2}A^T\nu$.

$q(\nu) = -\frac{1}{4}\nu^TAA^T\nu - \nu^Tb$.

Dual: $\max -\frac{1}{4}\nu^TAA^T\nu - \nu^Tb$.
\end{example}

\begin{example}[Maximum Entropy]
$\min \sum_i x_i \log x_i$ subject to $Ax = b$, $\mathbf{1}^Tx = 1$, $x \geq 0$.

Dual: $\max -\log\left(\sum_i \exp(-a_i^T\nu - \mu)\right)$ subject to dual
constraints, where $a_i$ is the $i$-th column of $A$.
\end{example}

\section{Min-Max Interpretation}

\begin{theorem}[Von Neumann--Sion Minimax Theorem]
If $\mathcal{X}$ is compact convex, $\mathcal{Y}$ is convex, and $\Phi(x,y)$ is
convex-concave (convex in $x$, concave in $y$), then:
\[
    \min_{x \in \mathcal{X}} \sup_{y \in \mathcal{Y}} \Phi(x,y)
    = \sup_{y \in \mathcal{Y}} \min_{x \in \mathcal{X}} \Phi(x,y).
\]
\end{theorem}

\begin{remark}
The Lagrangian $L(x, \lambda, \nu)$ is convex in $x$ (sum of convex functions) and
affine (hence concave) in $(\lambda, \nu)$. Under Slater's conditions, the minimax
theorem applies, yielding strong duality.
\end{remark}

\section{Saddle Point of the Lagrangian}

\begin{definition}[Saddle Point]
$(x^*, \lambda^*, \nu^*)$ is a \emph{saddle point} of the Lagrangian if:
\[
    L(x^*, \lambda, \nu) \leq L(x^*, \lambda^*, \nu^*) \leq L(x, \lambda^*, \nu^*),
    \quad \forall x, \; \forall \lambda \geq 0, \nu.
\]
\end{definition}

\begin{theorem}[Saddle Point and KKT Equivalence]
Under Slater's conditions, $(x^*, \lambda^*, \nu^*)$ is a saddle point of the
Lagrangian if and only if $x^*$ is primal optimal, $(\lambda^*, \nu^*)$ is dual
optimal, and strong duality holds.
\end{theorem}

\section{Economic Interpretation}

\begin{intuition}
Consider the profit maximization problem of a firm under resource constraints:
\[
    \max_x \; \mathrm{profit}(x) \quad \text{subject to} \quad
    \mathrm{resource}_i(x) \leq b_i.
\]
The multiplier $\lambda_i^*$ represents the \emph{shadow price} of resource $i$:
it is the marginal gain from increasing the availability of resource $i$ by one unit.
A competitive market would naturally set this price.
\end{intuition}

\section{Python Implementation}

\begin{codeblock}[title=\textbf{Lagrangian duality and duality gap}]
\begin{minted}{python}
import numpy as np
import cvxpy as cp

# Example: min ||x||^2 s.t. Ax = b
np.random.seed(42)
m, n = 3, 10
A = np.random.randn(m, n)
b = np.random.randn(m)

# Primal
x = cp.Variable(n)
primal = cp.Problem(cp.Minimize(cp.sum_squares(x)), [A @ x == b])
primal.solve()
p_star = primal.value
print(f"Primal value p* = {p_star:.6f}")
print(f"Multipliers: {primal.constraints[0].dual_value}")

# Analytical dual: max -1/4 nu^T A A^T nu - nu^T b
nu = cp.Variable(m)
M = A @ A.T
dual = cp.Problem(cp.Maximize(-0.25 * cp.quad_form(nu, M) - b @ nu))
dual.solve()
d_star = dual.value
print(f"Dual value d*   = {d_star:.6f}")
print(f"Duality gap     = {p_star - d_star:.2e}")

# LP: verify strong duality
c = np.random.randn(n)
A_ineq = np.random.randn(m, n)
b_ineq = np.abs(np.random.randn(m)) + 1

x_lp = cp.Variable(n)
primal_lp = cp.Problem(cp.Minimize(c @ x_lp),
                        [A_ineq @ x_lp <= b_ineq, x_lp >= 0])
primal_lp.solve()

y_lp = cp.Variable(m)
dual_lp = cp.Problem(cp.Maximize(b_ineq @ y_lp),
                      [A_ineq.T @ y_lp + c >= 0, y_lp <= 0])
dual_lp.solve()

print(f"\nLP primal: {primal_lp.value:.6f}")
print(f"LP dual:   {dual_lp.value:.6f}")
print(f"Gap:       {abs(primal_lp.value - dual_lp.value):.2e}")
\end{minted}
\end{codeblock}

\section{Exercises}

\begin{exercise}[$\star$]
Write the Lagrange dual of $\min c^Tx$ subject to $Ax = b$, $x \geq 0$.
\end{exercise}

\begin{exercise}[$\star$]
Compute the dual function for $\min \frac{1}{2}x^TQx + c^Tx$ subject to $Ax = b$
with $Q \succ 0$.
\end{exercise}

\begin{exercise}[$\star\star$]
Show that strong duality holds for linear programming without needing Slater's condition.
\end{exercise}

\begin{exercise}[$\star\star$]
Derive the dual of the soft-margin SVM problem and show it is a bounded QP.
\end{exercise}

\begin{exercise}[$\star\star$]
Show that $(x^*, \lambda^*, \nu^*)$ is a saddle point of the Lagrangian if and only
if the KKT conditions are satisfied.
\end{exercise}

\begin{exercise}[$\star\star\star$]
Prove Slater's theorem using the supporting hyperplane theorem applied to the set
$\mathcal{G}$ defined in the proof.
\end{exercise}

\begin{exercise}[$\star\star\star$]
Let $p^*(u)$ be the optimal value of the perturbed problem
$\min f(x)$ subject to $g_i(x) \leq u_i$. Show that $p^*$ is convex in $u$.
Deduce the interpretation of multipliers as subgradients of $p^*$.
\end{exercise}
